\begingroup
\small
\setlength{\tabcolsep}{4pt}
\renewcommand{\arraystretch}{1.12}
\begin{longtable}{|L{0.23\textwidth}|L{0.65\textwidth}|}
\caption{Спецификация прецедента «Поиск по естественно-языковому описанию»}\label{tab:spec-search}\\
\hline
\textbf{Раздел} & \textbf{Описание} \\
\hline
\endfirsthead
\multicolumn{2}{c}{Продолжение таблицы~\thetable}\\
\hline
\textbf{Раздел} & \textbf{Описание} \\
\hline
\endhead
Прецедент & Выполнить поиск аудиозаписей по естественно-языковому описанию \\
\hline
Краткое описание & Поиск аудиозаписей по текстовому описанию и дополнительным фильтрам с использованием семантического поиска по индексу. \\
\hline
Главные акторы & Пользователь. \\
\hline
Второстепенные акторы & Нет. \\
\hline
Предусловие & 1. Пользователь авторизован и имеет право выполнять поиск аудиозаписей.\newline 2. В поисковом индексе имеются данные об аудиозаписях. \\
\hline
Основной поток & 1. Пользователь вводит описание желаемой аудиозаписи.\newline 2. Пользователь при необходимости задаёт дополнительные фильтры по характеристикам аудиозаписей.\newline 3. Подсистема проверяет корректность описания и значений фильтров.\newline 4. Подсистема формирует векторное представление описания.\newline 5. Подсистема ищет аудиозаписи в поисковом индексе с учётом описания и фильтров.\newline 6. Подсистема упорядочивает найденные аудиозаписи по релевантности и предоставляет результаты пользователю. \\
\hline
Постусловия & 1. Пользователю предоставлен список подходящих аудиозаписей, упорядоченный по релевантности, либо сообщение об отсутствии результатов.\newline 2. Данные аудиозаписей и поискового индекса не изменены. \\
\hline
Альтернативные потоки & \textbf{E-1:} в описании или фильтрах есть ошибки. Подсистема сообщает пользователю об ошибке, поиск не выполняется.\newline\newline \textbf{E-2:} подходящие аудиозаписи не найдены. Подсистема сообщает пользователю об отсутствии результатов. \\
\hline
\end{longtable}
\endgroup
